\addcontentsline{toc}{chapter}{Executive Summary}
\chapter*{Executive Summary}

\begin{tcolorbox}
    This thesis examines the role of dedicated Data Institutions (DIs) in the design and development of machine learning (ML) datasets, emphasising their potential to `democratise AI'. The analysis critiques the current \textit{open AI} approach and its limitations, proposing that as an institutionalised form of commons-based data governance, DIs offer a structured democratic approach to data-centric AI policy-making. \\

    As an exploratory first pass, the research draws on multiple literatures, including critical political economy, data management, and governance theorising, to provide a normative and theoretical evaluation of DIs as an institutional template. This thesis proposes that integrating DIs into the ML value chain would leverage the strategic significance of datasets to create public value and insert a decentralised, community-accountable governance layer into the otherwise opaque ML development process. It sets out an agenda for DIs' collective governance mechanisms and fiduciary responsibilities to improve current (mal)practice in dataset development around data quality, usage, sustainability, and embedding of good governance practices.
\end{tcolorbox}